\clearpage
\section*{翻译说明}
\markboth{翻译说明}{翻译说明}
\phantomsection
\label{translation:notes}
\addcontentsline{toc}{section}{翻译说明}
本译稿依据 arXiv:1312.6114v11（2022 年 12 月 10 日）官方 TeX 源码及对应的 14 页 PDF，完整翻译正文、六个附录、图注、算法和原作者脚注，保留原文公式及对象编号。参考文献保留原始书目信息和引文标识；原图资产直接复用，图内英文标签保留。以下页码均指原 PDF。本节为译者新增，记录已核实并直接修入正文的问题，不表示已穷尽原文的所有错误。

\subsection*{已核实的公式与表述修正}
\begin{enumerate}
\item \textbf{第 3 页，第 2.2 节：梯度估计式。}
原式将梯度的下标写为分布 $q_{\bphi}(\bz)$，求和中的 $f$ 也没有代入样本。现统一为关于 $\bphi$ 的梯度，并使用 $f(\bzl)$，使采样分布与条件保持一致。依据为得分函数恒等式 $\nabla_{\bphi}q_{\bphi}=q_{\bphi}\nabla_{\bphi}\log q_{\bphi}$。正文明确 $f$ 不显含 $\bphi$；若显含，还须加入 $\mathbb{E}_{q_{\bphi}}[\nabla_{\bphi}f]$。

\item \textbf{第 3--4 页，式 (2)、(6)、(7) 及其说明。}
式 (2) 右端遗漏数据点上标 $i$，现与左端统一。式 (6) 的求和范围未用括号包住两个对数项，现补齐；式 (6)、(7) 的噪声采样上标统一为 $(i,l)$。第 4 页将第一项称为 KL 散度、将 $\log p$ 称为概率密度，现分别改为负 KL 散度和对数概率密度（或质量）。同时补全 minibatch 梯度中的 $\bphi$ 和估计量记号。

\item \textbf{第 4--5 页，第 2.4 节及第 3 节：重参数化。}
原文两处将 $g_{\bphi}$ 的输入顺序颠倒，现统一为 $g_{\bphi}(\beps,\bx)$。原证明直接对任意确定性映射写体积元等式，现改用推前分布下的期望恒等式，并把体积元等式限定于同维可逆可微情形，避免将可逆换元规则误用于一般映射。逆 CDF 方法中的均匀分布原写为 $\mathcal{U}(\bzero,\bI)$，现明确单变量取 $\mathcal{U}(0,1)$，独立分量逐一处理。

\item \textbf{第 5 页，第 2.4 节第 3 项：Gamma 与 Dirichlet 的构造。}
独立同率指数变量之和只得到正整数形状参数的 Gamma（埃尔朗）分布，不能代表任意形状参数的 Gamma，现补明限制。Dirichlet 不是 Gamma 变量的“加权和”，而是同尺度独立 Gamma 变量的归一化向量：若 $Y_j\sim\operatorname{Gamma}(\alpha_j,\beta)$ 独立，则 $Z_j=Y_j/\sum_kY_k$ 服从 $\operatorname{Dirichlet}(\bm\alpha)$。正文已按此修正。

\item \textbf{第 6 页，第 4 节：互信息与条件熵的符号。}
原文称最大化互信息等价于最大化条件熵，符号相反。固定输入分布时，$I(X;Z)=H(X)-H(X\mid Z)$，故应最大化负条件熵。又由条件 KL 散度非负，$\mathbb{E}\log p_{\bT}(X\mid Z)\leq-H(X\mid Z)$，所以期望对数似然下界对应的是负条件熵。正文已修正这两处符号，并将朴素梯度估计量的节号从 2.1 改为实际所在的 2.2。

\item \textbf{第 6 页，第 5 节：编码器与解码器。}
原文把生成模型标为 encoder，把变分近似标为 decoder，与第 2.1 节定义颠倒。现分别改为解码器和编码器；实验设置和数据不变。

\item \textbf{第 10--11 页，附录 B 和 C。}
附录 B 推导将近似分布误写为 $q_{\bT}$，现统一为 $q_{\bphi}$，并删去 KL 记号中多余的左括号。式 (11) 的伯努利对数似然补上求和括号，使 $x_i\log y_i$ 和 $(1-x_i)\log(1-y_i)$ 两项均对 $i$ 求和。

\item \textbf{第 12 页，附录 D：边缘似然估计量。}
两处样本求和的分母仍写 $p_{\bT}(\bz)$，现改为 $p_{\bT}(\bzl)$。补明 $q$ 归一化且支撑集位于联合密度为正的区域：此时 $\mathbb{E}_{p_{\bT}(\bz\mid\bxi)}[q(\bz)/p_{\bT}(\bxi,\bz)]=1/p_{\bT}(\bxi)$。有限样本均值取倒数一般有偏，正文保留“估计”而不增称无偏。

\item \textbf{第 13 页，附录 F：分布名、条件与引用。}
式 (16)、(18) 的近似后验条件统一为 $\bxi$；说明解析 KL 的高斯条件中，原 $p_{\bT}(\bx)$ 改为潜变量先验 $p_{\bT}(\bz)$，并明确全局参数情形的对应分布。重参数化条件的残缺节号补为 2.4；随机梯度算法引用由算法 1 改为对应的算法 2。例子中的参数先验 $p_{\balpha}(\bT)=\mathcal{N}(\bz;\bzero,\bI)$ 改为 $\mathcal{N}(\bT;\bzero,\bI)$。

\item \textbf{第 13 页，式 (22)：两层下界的区别与数据采样。}
原式把替换潜变量对数边缘似然后的联合下界仍记成式 (14) 的 $\mathcal{L}(\bphi;\bX)$，同时只说给定单个数据点，却未交代式中的 $\bxl$ 如何产生。现以 $\mathcal{J}$ 区分联合下界，并明确 $\bxl$ 从数据集独立均匀抽样。依据为
\[
\mathcal{L}(\bphi;\bX)-\mathcal{J}(\bphi;\bX)
=\sum_{i=1}^N\mathbb{E}_{q_{\bphi}(\bT)}
\left[D_{KL}\bigl(q_{\bphi}(\bz\mid\bxi)\Vert p_{\bT}(\bz\mid\bxi)\bigr)\right]\geq0.
\]
式 (21) 的 $N$ 倍局部项经均匀抽样后的期望正好恢复对全部数据点的求和；只有潜变量近似后验精确时，两层下界才相等。

\item \textbf{第 14 页，式 (24)：完全 VB 的高斯估计量。}
原重构项写为 $\log p_{\bT}(\bxi\bzi)$，缺少条件符号，数据及潜变量索引与外层求和不一致，也没有明确参数样本。现改为 $\log p_{\bT^{(l)}}(\bxl\mid\bzl)$，并给出相应重参数化样本。全局参数 KL 项原使用潜变量维数 $J$ 且保留游离的样本索引 $l$，现改为参数维数 $P$，去掉不属于全局变分参数的样本索引。左端同步改为联合下界 $\mathcal{J}$。这来自“期望重构对数似然减去局部 KL 与全局 KL”的分解。
\end{enumerate}

\subsection*{记号与排版统一}
数学粗体统一使用 \texttt{\textbackslash bm}，微分算子统一为正体 $\mathrm{d}$；若使用迹运算符，采用小写正体 $\operatorname{tr}$。原文以方差向量乘单位矩阵简记对角协方差，现显式写为 $\operatorname{diag}(\bsigma^2)$，其中平方按分量计算，不改变分布含义。噪声样本 $\bzeta^{(l)}$ 的向量字形保持一致。原文各处“边缘似然”与“对数边缘似然”按公式实际含义区分。

采用既有蓝色标题中文单栏模板，分页因此与原文不同。原始英文 PDF 和 TeX 源包独立保存，修正仅作用于译稿。书目中的历史元数据按原稿保留；本次未重新验证实验结果，也不把排版统一当作作者勘误。
